\section{座標と教師の契約}
\label{app:contracts}

\subsection{変換の向きと姿勢の直列化}

保存する剛体変換はすべて行優先の \(4\times4\) 同次行列であり、
列ベクトルへの作用は
\begin{equation}
  \begin{bmatrix}\mathbf{x}_b\\1\end{bmatrix}
    =\T{b}{a}\begin{bmatrix}\mathbf{x}_a\\1\end{bmatrix},
  \qquad
  \T{a}{b}=\bigl(\T{b}{a}\bigr)^{-1}
\end{equation}
である。回転は正規直交かつ行列式 \(+1\) を要求する。
一様尺度はメートルアダプタが持ち、
名目上の剛体変換の中へ隠さない。
生成サンプルのカメラ変換はカメラからシーンであり、
世界からカメラへの投影は常にその逆変換を使う。

データセットは生のカメラ・コート変換と内部パラメータを保存し、
補正に追随する消費者側が
\begin{center}
\texttt{tx, ty, tz, a11, a12, a13, a21, a22, a23, logf}
\end{center}
の順で姿勢ベクトルを導出する。
並進は正準コート座標でのカメラ中心、
六つの回転値はカメラから正準コートへの回転の最初の二行、
焦点距離は補正後モデル画素単位である。
この直列化では四元数の順序規約が問題にならない。

\subsection{十四のキーポイントチャネル}

明示的に選択したデータセットは、
カメラ相対の意味チャネルと不変の物理番号の両方を保持する。
カメラがコートの負の \(Y\) 側にあるときを同一視なしの対応、
正の \(Y\) 側にあるときを半回転の対応とする。
どちらの場合も、
意味チャネルは十四個の物理点の置換として与えられる。

カメラ中心が中央面から \(10^{-6}\) m 以内にある場合は、
近い側と遠い側が幾何的に定義できないため、
レンダリングの前に明示的に棄却する。
左右の割り当ては常にコート局所の \(X\) に従う。

\subsection{可視性の層}

四つの状態を区別して保存する。

\begin{enumerate}[leftmargin=1.6em]
  \item \emph{カメラ前方}：カメラ深度が正である。
  \item \emph{画面内}：投影座標が画像の内側にある。
  \item \emph{描画支持}：半径 1 画素の近傍に十分な alpha と
  正の描画深度がある（alpha の閾値は \(0.01\)）。
  これは点単位の深度遮蔽判定ではない。
  \item \emph{教師有効}：上の生成物の論理積。
\end{enumerate}

教師有効の定義は生成側の権威であり、
予測スコアや予測深度で事後に変えない。
補正は画像フレームの幾何と教師有効性を
同じ手順で再計算する。
バッチ化で生じる詰め物の画布は、
決して有効な内容として扱わない。

\subsection{データセットの項目}

各サンプルは、
メートル系シーン座標でのカメラ変換、
シーンから対象コートへの変換、
補正前の内部パラメータ、
物理番号付きのキーポイント教師、
領域と線の教師、
および可視性の層を保持する。
姿勢ベクトルは保存せず、
補正に追随するアダプタが導出する。
